% ===========================================================================
\chapter{Flattening the Earth}\label{ch:projection}

A screen is flat. To draw a map we need every position as an $x$ and a $y$
on a flat surface, and it would be nice if those were in metres, so that
the ordinary geometry of the flat plane (Pythagoras' theorem, straight
lines, areas of polygons) just works. Turning latitude and longitude into
flat coordinates is called a \emph{map projection}.

\section{You can't flatten an orange peel}

Try to flatten a piece of orange peel and it tears or stretches. A sphere
cannot be laid flat without distorting it somehow, and every projection is
a choice about what to distort. Some keep areas right and bend shapes; some
keep shapes right in small patches and blow up areas far from the equator;
none keeps everything right everywhere.

\begin{aside}
\textbf{Web maps and Mercator.} Almost every map on the web, from
OpenStreetMap to the big commercial ones, uses the \emph{Web Mercator}
projection:
\[
  x = R\lambda, \qquad y = R \ln\tan\Big(\frac{\pi}{4} + \frac{\varphi}{2}\Big).
\]
Its great virtue is that it keeps angles right: a street corner that is a
right angle on the ground is a right angle on the screen, at every zoom.
Its great vice is scale: a metre on the ground is stretched by
$1/\cos\varphi$ on the map, more and more towards the poles. That is why
Greenland looks as big as Africa on a web map, when it is fourteen times
smaller. A map app that shows Mercator still has to measure distances on
the sphere, or correct for the stretch, as we do.
\end{aside}

The good news is that the smaller the area, the less any of this matters. A
piece of orange peel the size of a fingernail lies almost flat. Kestrel
Island is about four kilometres across, a fingernail on a twelve-thousand
kilometre ball, and the simplest projection there is will be accurate to a
few centimetres.

\section{How a degree of longitude shrinks}

The parallel at latitude $\varphi$ is a circle. Slice the Earth through
the poles and look at it side on:

\begin{center}
\begin{tikzpicture}[scale=1.2]
  \draw[accent!60, fill=titlebg!40] (0,0) circle (2);
  \draw[color=cm, dashed] (0,-2.3) -- (0,2.3) node[above, font=\small] {axis};
  \draw[color=cm] (-2.2,0) -- (2.2,0) node[right, font=\small] {equator};
  \coordinate (P) at (47:2);
  \draw[st, thick] (0,0) -- (P) node[midway, above left, font=\small] {$R$};
  \draw[accent, thick] (0,{2*sin(47)}) -- (P) node[midway, above, font=\small] {$r = R\cos\varphi$};
  \draw[color=cm, dashed] ({-2*cos(47)},{2*sin(47)}) -- (0,{2*sin(47)});
  \draw (0.55,0) arc[start angle=0, end angle=47, radius=0.55];
  \node[font=\small] at (24:0.85) {$\varphi$};
  \fill[st] (P) circle (0.06);
\end{tikzpicture}
\end{center}

The parallel's radius is the side of a right-angled triangle whose
hypotenuse is the Earth's radius $R$, with the angle $\varphi$ at the
centre, so it is $r = R\cos\varphi$. A degree of longitude is $1/360$ of
that circle, $2\pi R \cos\varphi / 360$: the 111,195 metres of a degree at
the equator, times $\cos\varphi$. At Kestrel Island's latitude,
$\cos 47^\circ \approx 0.68$, so a degree east is only about 75.8 km, while
a degree north is still 111.2 km.

\section{A local projection}

So here is the projection. Pick a centre for the map, at latitude
$\varphi_0$ and longitude $\lambda_0$. Then a place at $\varphi$, $\lambda$
(in degrees) goes to
\[
  x = (\lambda - \lambda_0) \cdot k_\lambda, \qquad
  y = (\varphi - \varphi_0) \cdot k_\varphi,
\]
where $k_\varphi = 2\pi R/360$ metres per degree of latitude, and
$k_\lambda = k_\varphi \cos\varphi_0$ metres per degree of longitude,
measured at the centre. The centre goes to $(0, 0)$; $x$ grows to the east
and $y$ to the north, both in metres. This is called the
\emph{equirectangular} projection.

It is exactly right along the centre's parallel, where we measured
$k_\lambda$. North or south of it, the true size of a degree of longitude
differs a little from $k_\lambda$. How little? Moving 1.6 km north, the
distance from the centre to Kestrel Island's northern tip, changes the
latitude by $0.0144^\circ$, which changes $\cos\varphi$ by about 0.03\%.
So an east-west distance up there comes out 0.03\% too long: 30
centimetres in a kilometre. For a road map, that is perfect.

Going back is just as easy: $\lambda = \lambda_0 + x / k_\lambda$ and
$\varphi = \varphi_0 + y / k_\varphi$.

\section{Points on the flat map}

A flat position is a \code{Point} with an \code{x} and a \code{y}. We will
use points in two ways: as positions, and as \emph{arrows} (vectors), such
as ``the arrow from the harbour to the lighthouse'', which is
\code{lighthouse - harbour}. Adding an arrow to a position moves the
position along the arrow; multiplying an arrow by a number stretches it.

Nex lets a class give a method the name of an arithmetic operator with
\code{alias}, so that \code{a + b} on two points means \code{a.plus(b)}.
That keeps geometry readable: \code{a + (b - a) * 0.5} is the point halfway
from \code{a} to \code{b}.

\begin{append}{geo.nex}
-- A position on a flat map, in metres: x east and y north of the map's
-- centre.
class Point
create
  make(px, py: Real) do
    x := px
    y := py
  end
feature
  once x: Real
  once y: Real

  plus(other: Point): Point alias "+" do
    result := create Point.make(x + other.x, y + other.y)
  end

  minus(other: Point): Point alias "-" do
    result := create Point.make(x - other.x, y - other.y)
  end

  -- This point scaled by k, as an arrow from the origin.
  times(k: Real): Point alias "*" do
    result := create Point.make(x * k, y * k)
  end
\end{append}

Two more operations on arrows. The \emph{length} of an arrow $(x, y)$ is
$\sqrt{x^2 + y^2}$, by Pythagoras; Nex's \code{hypot} computes exactly
that. The \emph{dot product} of two arrows, $a_x b_x + a_y b_y$, will be
useful in the next chapter, where we explain what it measures. The distance
between two points is the length of the arrow from one to the other.

\begin{continued}{geo.nex}
  -- The dot product, treating both points as arrows from the origin.
  dot(other: Point): Real do
    result := x * other.x + y * other.y
  end

  -- The length of this point as an arrow from the origin.
  length(): Real do
    result := x.hypot(y)
  end

  distance_to(other: Point): Real do
    result := (other - this).length
  ensure
    non_negative: result >= 0.0
  end

  -- "(123.4, -56.7)"
  to_string(): String do
    result := "(" + fixed(x, 1) + ", " + fixed(y, 1) + ")"
  end
end
\end{continued}

\section{The projection}

The projection holds the centre and the two scale factors, worked out once:

\begin{append}{geo.nex}
-- The local flat projection of the area around `centre`: a small piece of
-- the globe laid flat, with `centre` at (0, 0), x pointing east and y
-- north, both in metres.
--
-- A degree of latitude is the same length everywhere, but a degree of
-- longitude shrinks towards the poles, in proportion to the cosine of the
-- latitude: meridians meet at the poles. This projection measures both at
-- the centre's latitude. That is exact along the centre's parallel and
-- very nearly exact a few kilometres away, which is all a town or an
-- island needs. (A whole country or continent needs a cleverer projection.)
class Projection
create
  make(c: Lat_Lon) do
    centre := c
    metres_per_degree_lat := Earth.RADIUS * Real.PI / 180.0
    metres_per_degree_lon := metres_per_degree_lat * c.latitude.to_radians.cos
  end
feature
  centre: Lat_Lon
  metres_per_degree_lat: Real
  metres_per_degree_lon: Real

  to_point(p: Lat_Lon): Point do
    result := create Point.make(
      (p.longitude - centre.longitude) * metres_per_degree_lon,
      (p.latitude - centre.latitude) * metres_per_degree_lat)
  end

  to_lat_lon(p: Point): Lat_Lon do
    result := create Lat_Lon.make(
      centre.latitude + p.y / metres_per_degree_lat,
      centre.longitude + p.x / metres_per_degree_lon)
  ensure
    round_trip: to_point(result).distance_to(p) < 0.001
  end
invariant
  positive_scales: metres_per_degree_lat > 0.0 and metres_per_degree_lon > 0.0
end
\end{append}

The postcondition of \code{to_lat_lon} checks the two directions against
each other: projecting the answer forward again must land within a
millimetre of where we started. A contract like this is cheap insurance.
If someone later ``improves'' one direction and forgets the other, the
program stops at the first mismatch instead of drawing roads a few metres
off.

\begin{tryit}{try_projection.nex}
intern Geo

let harbour: Lat_Lon := create Lat_Lon.make(47.052428, -34.569106)
let lighthouse: Lat_Lon := create Lat_Lon.make(47.079318, -34.565277)
let flat: Projection := create Projection.make(create Lat_Lon.make(47.066, -34.5695))
let h: Point := flat.to_point(harbour)
let l: Point := flat.to_point(lighthouse)
say("harbour    " + h.to_string())
say("lighthouse " + l.to_string())
say("on the sphere:   " + fixed(harbour.distance_to(lighthouse), 2) + " m")
say("on the flat map: " + fixed(h.distance_to(l), 2) + " m")
say("back again: " + flat.to_lat_lon(l).to_string())

-- Pretending a degree is the same length both ways:
let k: Real := flat.metres_per_degree_lat
let naive: Point := create Point.make((lighthouse.longitude - harbour.longitude) * k,
  (lighthouse.latitude - harbour.latitude) * k)
say("degrees as if they were squares: " + fixed(naive.length(), 2) + " m")
\end{tryit}
\tryoutput{try_projection}

The flat map agrees with the sphere to within a few centimetres over three
kilometres. Pretending that a degree of longitude is as long as a degree of
latitude gets the distance wrong by several metres here, and it would be
much worse for two places side by side east and west. A map drawn that way
would look stretched sideways, by $1/\cos 47^\circ$, nearly one and a
half times.
